\section{实验分析与总结}  % 原来是 \chapter

\subsection{网络性能分析}

\begin{itemize}
\item 三层网络已能较好拟合 MNIST
\item ReLU 提高非线性表达能力
\item Softmax + 交叉熵稳定训练
\end{itemize}

\subsection{参数影响}

\begin{table}[htbp]
\centering
\begin{tabular}{|l|l|}
\hline
\textbf{参数} & \textbf{影响} \\
\hline
层数增加 & 表达能力增强但易过拟合 \\
神经元数量 & 过少欠拟合，过多计算量大 \\
学习率 & 过大震荡，过小收敛慢 \\
batch\_size & 影响稳定性与速度 \\
\hline
\end{tabular}
\caption{参数影响分析}
\label{tab:impact}
\end{table}

\subsection{优化效果}

\begin{itemize}
\item 动量：加快收敛
\item L2正则化：防止过拟合
\item 学习率衰减：提升后期稳定性
\end{itemize}

\subsection{实验总结}

本实验通过纯 NumPy 实现了一个完整的神经网络训练流程，深入理解了：

\begin{itemize}
\item 神经网络前向传播与反向传播机制
\item 各种层（Matmul、ReLU、Softmax）的作用
\item 损失函数与优化算法原理
\item 深度学习训练流程
\end{itemize}

该实验强化了对深度学习底层实现的理解，为后续使用框架（如 PyTorch）奠定了基础。